\documentclass[12pt]{article}

% Packages
\usepackage{amsmath, amssymb, amsthm}
\usepackage{geometry}
\usepackage{fancyhdr}
\usepackage{enumitem}
\usepackage{titlesec}
\usepackage{tcolorbox}

% Page Setup
\geometry{margin=1in}
\setlength{\parindent}{0pt}
\setlength{\parskip}{0em}
\setcounter{secnumdepth}{3}

\pagestyle{fancy}
\fancyhf{}
\lhead{Ethan Fan}
\rhead{PCH}
\cfoot{\thepage}

\allowdisplaybreaks

% Section Formatting
\titleformat{\section}
  {\bfseries}
  {}
  {0em}
  {}
\titlespacing*{\section}{0pt}{0pt}{0.3em}

\titleformat{\subsection}[runin]
  {\bfseries}
  {}
  {0em}
  {}
\titlespacing*{\subsection}{0pt}{0pt}{0em}

\titleformat{\subsubsection}[runin]
  {\itshape}
  {(\roman{subsubsection})}
  {0em}
  {}

% mathematical commands
\newcommand{\ddt}{\frac{d}{dt}}
\newcommand{\ddx}{\frac{d}{dx}}
\newcommand{\dydx}{\frac{dy}{dx}}

\newcommand{\definition}[1]{\begin{tcolorbox}[colback=red!5!white,colframe=red!75!black,parbox=false] #1 \end{tcolorbox}}

\newcommand{\theorem}[2]{\begin{tcolorbox}[title={#1},colback=blue!5!white,colframe=blue!75!black,parbox=false] #2 \end{tcolorbox}}

\newcommand{\example}[2]{\begin{tcolorbox}[title={Example: #1},colback=brown!5!white,colframe=brown!75!black,parbox=false] #2 \end{tcolorbox}}

\newcommand{\remark}[2]{\begin{tcolorbox}[title={#1},colback=black!5!white,colframe=black!75!black,parbox=false] #2 \end{tcolorbox}}

% Document
\begin{document}

\vspace*{0cm}

\begin{center}
    {\Large \bfseries Problem Set 58} \\[0.5cm]
    {\large The Squeeze Theorem} \\[0.35cm]
    {\normalsize March 24, 2026}
\end{center}

% Questions

\section{Question 3}
Assume that $\lim\limits_{\theta \to -1^-}f(\theta)$ exists and $\dfrac{\theta^2+\theta-2}{\theta +3} \le \dfrac{f(\theta)}{\theta^2} \le \dfrac{\theta^2+2\theta-1}{\theta +3}$. Find $\lim\limits_{\theta \to -1^-}f(\theta)$.
\begin{gather*}
    \dfrac{\theta^2+\theta-2}{\theta +3} \le \dfrac{f(\theta)}{\theta^2} \le \dfrac{\theta^2+2\theta-1}{\theta +3}\\
    \dfrac{\theta^2(\theta^2+\theta-2)}{\theta +3} \le f(\theta) \le \dfrac{\theta^2(\theta^2+2\theta-1)}{\theta +3}\\
    \lim\limits_{\theta \to -1^-}\dfrac{\theta^2(\theta^2+\theta-2)}{\theta +3}=\lim\limits_{\theta \to -1^-}\frac{1(1-1-2)}{2}=-1\\
    \lim\limits_{\theta \to -1^-}\dfrac{\theta^2(\theta^2+2\theta-1)}{\theta +3}=\lim\limits_{\theta \to -1^-}\frac{1(1-2-1)}{2}=-1\\
    \lim\limits_{\theta \to -1^-}\dfrac{\theta^2(\theta^2+\theta-2)}{\theta +3}=\lim\limits_{\theta \to -1^-}\dfrac{\theta^2(\theta^2+2\theta-1)}{\theta +3}=-1\\
    \boxed{\lim\limits_{\theta \to -1^-}f(\theta)=-1}
\end{gather*}

\section{Question 5}
Let $m$ and $n$ be positive integers greater than 1. Find the value of $\dfrac{m}{n}$ if
\[
\lim\limits_{\alpha \to 0} \frac{e^{\cos (\alpha ^n)}-e}{\alpha ^ m}=-\frac{e}{2}
\]
\begin{gather*}
    \lim\limits_{\alpha \to 0} \frac{e^{(1-\frac{\alpha ^ {2n}}{2})}-e}{\alpha ^ m}=-\frac{e}{2}\\
    \lim\limits_{\alpha \to 0} \frac{e^{(-\frac{\alpha ^ {2n}}{2})}-1}{\alpha ^ m}=-\frac{1}{2}\\
    \lim\limits_{\alpha \to 0} (\frac{e^{(-\frac{\alpha ^ {2n}}{2})}-1}{-\frac{\alpha ^ {2n}}{2}})(\frac{-\frac{\alpha ^ {2n}}{2}}{\alpha ^m})=-\frac{1}{2}\\
    \lim\limits_{\alpha \to 0}\frac{-\frac{\alpha ^ {2n}}{2}}{\alpha ^m}=-\frac{1}{2}\\
    -\frac{1}{2} \lim\limits_{\alpha \to 0} \alpha ^{2n-m}=-\frac{1}{2} \\
    \lim\limits_{\alpha \to 0} \alpha ^{2n-m}=1\implies 2n-m=0\\
    \boxed{\frac{m}{n}=2}
\end{gather*}

\section{Question 6}
Find the value of $\theta$ such that
\[
\lim\limits_{x \to 0} (1+x\ln (1+b^2))^{\frac{1}{x}}=2b\sin ^2 \theta, \quad b>0, \quad \theta \in (-\pi,\pi].
\]
\begin{gather*}
    \lim\limits_{x \to 0} (1+x\ln (1+b^2))^{\frac{ln(1+b^2)}{xln(1+b^2)}}=2b\sin ^2 \theta\\
    e ^{ln(1+b^2)}=2b\sin ^2 \theta\\
    \sin ^2 \theta = \frac{1+b^2}{2b}\\
    \frac{1+b^2}{2b}\ge 2\sqrt{\frac{1}{2b}\cdot \frac{b^2}{2b}}=1\\
    \sin ^2 \theta = 1\\
    \boxed{\theta = \frac{\pi}{2} \text{ or } -\frac{\pi}{2}}
\end{gather*}

\section{Question 7}
\subsection{(a)} \, Find the perimeter of a regular $n$-gon inscribed in a circle of radius $r$.
\begin{gather*}
    s^2=r^2+r^2-2r^2\cos (\frac{2\pi}{n})\\
    s=r\sqrt{2(1-\cos (\frac{2\pi}{n}))}\\
    s=r\sqrt{2(1-1+2 \sin^2 \frac{\pi}{n})}\\
    s=2r \sin \frac{\pi}{n}\\
    \boxed{P=2nr \sin \frac{\pi}{n}}
\end{gather*}

\subsection{(b)} \, What value does this perimeter approach as $n$ becomes very large?
\begin{gather*}
    \lim_{n \to \infty}2nr \sin \frac{\pi}{n}=2\pi r\lim_{n \to \infty} \frac{\sin \frac{\pi}{n}}{\frac{\pi}{n}}=2\pi r\\
    \boxed{\lim_{n \to \infty}P=2\pi r}
\end{gather*}

\subsection{(c)} \, What limit can you guess from this?
\begin{gather*}
    \boxed{\lim_{x \to 0}\frac{\sin x}{x}=1}
\end{gather*}

\end{document}

